\documentclass[a4paper,12pt]{article}
\usepackage[T2A]{fontenc}
\usepackage[utf8]{inputenc}
\usepackage[russian]{babel}
\usepackage{amsmath,amssymb}
\usepackage{geometry}
\geometry{left=24mm,right=24mm,top=20mm,bottom=23mm}
\usepackage{xcolor}
\usepackage{microtype}
\usepackage{booktabs,tabularx}
\usepackage{enumitem}
\usepackage{fancyhdr}
\usepackage{tikz}
\usepackage{hyperref}
\hypersetup{hidelinks}
\definecolor{accent}{HTML}{087F83}
\definecolor{ink}{HTML}{26353A}
\definecolor{muted}{HTML}{5D6B71}
\color{ink}
\setlength{\parindent}{0pt}
\setlength{\parskip}{5pt}
\linespread{1.09}
\setlist[enumerate]{leftmargin=*,itemsep=4pt,topsep=3pt}
\pagestyle{fancy}
\fancyhf{}
\renewcommand{\headrulewidth}{0pt}
\renewcommand{\footrulewidth}{0pt}
\fancyfoot[C]{\footnotesize\color{muted}Лёвин Артём Александрович эксклюзивно для Ярослава Верещагина}
\fancyfoot[R]{\footnotesize\color{muted}\thepage}
\newcommand{\topic}[1]{\par\medskip{\large\bfseries\color{accent}#1}\par\smallskip}
\newcommand{\eqtask}[1]{\begin{center}\vspace{2mm}\Large $#1$\vspace{2mm}\end{center}}
\begin{document}
\thispagestyle{empty}
\begin{center}
\vspace*{23mm}
{\large\color{accent}МАТЕМАТИКА}\par
\vspace{18mm}
{\LARGE Ревью домашней работы}\par
\vspace{12mm}
{\Large\color{accent}Метод интервалов\par
Рациональные неравенства}\par
\vspace{20mm}
{\normalsize Разбор заданий 4 и 5}\par
\vfill
{\large Ярослав Верещагин}\par
\vspace{10mm}
{29 октября 2026 г.}\par
\vspace{16mm}
{\small Лёвин Артём Александрович}\par
{\small эксклюзивно для Ярослава Верещагина}\par
\vspace{19mm}
\begin{tikzpicture}[x=1cm,y=1cm]
\draw[accent,line width=0.9pt] (-4.5,0) .. controls (-3.3,0.55) and (-2.2,-0.35) .. (-1.1,0)
.. controls (0.15,0.48) and (1.1,-0.35) .. (2.2,0)
.. controls (3.0,0.27) and (3.8,0.10) .. (4.5,0);
\end{tikzpicture}
\vspace{5mm}
\end{center}
\clearpage

{\footnotesize\color{muted}29.10.2026 \hfill МЕТОД ИНТЕРВАЛОВ}
\vspace{4mm}

{\LARGE\bfseries\color{accent}Задание 4}\par
\eqtask{\frac{(x-1)^2}{x+2}\ge 0}

\topic{Где возникла ошибка}
В Вашей схеме знак дроби меняется при $x=1$. Однако множитель $(x-1)^2$ имеет \textbf{чётную степень}: слева и справа от единицы он положителен, а в самой точке равен нулю. Поэтому при переходе через $x=1$ знак выражения не меняется.

В записанном ответе точка $x=1$ исключена, хотя условие допускает равенство нулю. Кроме того, при подстановке $x=0$ в проверочное выражение был пропущен квадрат.

\topic{Правильное решение}
\begin{enumerate}
\item \textbf{Область допустимых значений.} Знаменатель не может быть равен нулю:
\[
x+2\ne0\quad\Longrightarrow\quad x\ne-2.
\]
\item \textbf{Числитель.} При любом $x$ квадрат неотрицателен:
\[
(x-1)^2\ge0,\qquad (x-1)^2=0\ \text{только при}\ x=1.
\]
\item \textbf{Знаки на промежутках.} Если $x\ne1$, числитель положителен. Следовательно, знак дроби определяется знаменателем $x+2$:
\[\renewcommand{\arraystretch}{1.4}
\begin{array}{ccccc}
\toprule
x&(-\infty,-2)&(-2,1)&x=1&(1,+\infty)\\
\midrule
\dfrac{(x-1)^2}{x+2}&-&+&0&+\\
\bottomrule
\end{array}
\]
\item \textbf{Проверка особой точки.} При $x=1$ дробь определена и равна нулю:
\[
\frac{(1-1)^2}{1+2}=0.
\]
При $x=0$ получаем $\dfrac{(0-1)^2}{0+2}=\dfrac12>0$.
\end{enumerate}

\textbf{Ответ:}\quad $\boxed{x\in(-2,+\infty)}$.

\topic{Задача на закрепление}
Решите методом интервалов, отдельно проверив точку, в которой квадрат обращается в нуль:
\[
\frac{(x+1)^2}{x-3}\ge0.
\]
\clearpage

{\footnotesize\color{muted}29.10.2026 \hfill МЕТОД ИНТЕРВАЛОВ}
\vspace{4mm}

{\LARGE\bfseries\color{accent}Задание 5}\par
\eqtask{\frac{(x-4)^2(x+1)}{x-2}>0}

\topic{Где возникла ошибка}
На числовой прямой левее точки $-1$ Вы правильно поставили знак «плюс», но при записи ответа \textbf{потеряли весь этот промежуток}. Кроме того, в проверке значения при $x=0$ квадрат $(x-4)^2$ был записан как обычный множитель $x-4$: такая замена меняет знак и делает проверку некорректной.

\topic{Правильное решение}
\begin{enumerate}
\item \textbf{ОДЗ.} Знаменатель $x-2$ не равен нулю, поэтому $x\ne2$.
\item \textbf{Особые точки.} Числитель равен нулю при $x=-1$ и $x=4$. В точке $x=4$ множитель $(x-4)^2$ имеет чётную степень, поэтому знак при переходе через неё сохраняется.
\item \textbf{Определяем знак.} Вне точки $x=4$ квадрат положителен. Достаточно определить знак дроби $\dfrac{x+1}{x-2}$:
\[\renewcommand{\arraystretch}{1.4}
\begin{array}{ccccc}
\toprule
x&(-\infty,-1)&(-1,2)&(2,4)&(4,+\infty)\\
\midrule
\dfrac{(x-4)^2(x+1)}{x-2}&+&-&+&+\\
\bottomrule
\end{array}
\]
\item \textbf{Проверка пропущенного промежутка.} Возьмём $x=-2$:
\[
\frac{(-2-4)^2(-2+1)}{-2-2}
=\frac{36\cdot(-1)}{-4}=9>0.
\]
Значит, все $x<-1$ подходят.
\item \textbf{Границы.} Точки $x=-1$ и $x=4$ дают нуль и исключаются, поскольку неравенство строгое. Точка $x=2$ запрещена по ОДЗ.
\end{enumerate}

\textbf{Ответ:}
\[
\boxed{x\in(-\infty,-1)\cup(2,4)\cup(4,+\infty)}.
\]

\topic{Задача на закрепление}
Решите методом интервалов и внимательно перенесите все подходящие промежутки в итоговый ответ:
\[
\frac{(x+2)^2(x-3)}{x+1}>0.
\]
\smallskip
{\footnotesize\color{muted}\textbf{Главное правило.} Множитель в чётной степени не меняет знак при переходе через свой нуль. Сам нуль включается в ответ при нестрогом знаке и допустимом знаменателе; при строгом знаке он исключается.}

\end{document}
